\paragraph{Proofs - Do-Calculus}\hfill%


\begin{Lem}
    \label{lem:int-null-set}
    For pairwise disjoint $B,C \ins V$ and $D \ins V \cup J$ and a measurable subset $N \ins \Xcal_{B \cup C\cup D \cup J}$ the following statements are equivalent:
    \begin{enumerate}
        \item $N$ is a $P(X_B,X_C|\doit(X_{I_B},X_{D\cup J}))$-null set.
        \item For every decomposition $B=B_1 \dcup B_2$ the set $N$ is a 
            $P(X_B,X_C|\doit(X_{B_2},X_{D\cup J}))$-null set.
        \item For every decomposition $B=B_1 \dcup B_2$ the set $N$ is a 
            $P(X_{B_1},X_C|\doit(X_{B_2},X_{D\cup J}))$-null set.
    \end{enumerate}
\begin{proof}
    Every value $x_{I_B}=(x_{I_v})_{v \in B} \in \Xcal_{I_B}$ defines a decomposition $B=B_1 \dcup B_2$ via:
    \[ B_1 := \lC v \in B \st x_{I_v} = \star \rC, \qquad B_2 := \lC v \in B \st x_{I_v} \in \Xcal_v \rC.  \]
    So running through all values $x_{I_B} \in \Xcal_{I_B}$ is the same as running through all subsets $B_2 \ins B$ and all values $x_{B_2} \in \Xcal_{B_2}$, while putting $x_{I_{B_1}}=\star$ for $B_1 = B \sm B_2$.
    Furthermore, we have the following identities:
    \begin{align*}
          &   P((X_B,X_C) \in N_{x_{D\cup J}}|\doit(X_{I_B}=x_{I_B},X_{D\cup J}=x_{D\cup J})) \\
        &=   P((X_B,X_C) \in N_{x_{D\cup J}}|\doit(X_{I_{B_1}} = \star ,X_{I_{B_2}}=x_{B_2},X_{D\cup J}=x_{D\cup J})) \\
     &=   P((X_B,X_C) \in N_{x_{D\cup J}}|\doit(X_{B_2}=x_{B_2},X_{D\cup J}=x_{D\cup J})) \\
     &= \lp P(X_{B_1},X_C|\doit(X_{B_2}=x_{B_2},X_{D\cup J}=x_{D\cup J})) \otimes \delta(X_{B_2}|X_{B_2}=x_{B_2})\rp\lp N_{x_{D\cup J}} \rp \\
     &= P((X_{B_1},X_C) \in N_{(x_{B_2},x_{D\cup J})}|\doit(X_{B_2}=x_{B_2},X_{D\cup J}=x_{D\cup J})).
    \end{align*}
    So the first line vanishes for all values $x_{I_B} \in \Xcal_{I_B}$ and $x_{D\cup J} \in \Xcal_{D\cup J}$ if and only if any other line vanishes
    for all subsets $B_2 \ins B$ and all values $x_{B_2} \in \Xcal_{B_2}$ and $x_{D\cup J} \in \Xcal_{D\cup J}$.
    This shows the claim.
\end{proof}
\end{Lem}

%\Joris{Suggestion: explain what it means that a Markov kernel $K(X_A|X_B,X_C,X_D)$ is unique up to a measurable $K(X_B|X_C,X_D)$-null set in $\Xcal_{B\cup C}$ (rather than in $\Xcal_{B \cup C \cup D}$): that the null set can be considered constant in $X_{D \sm (B \cup C)}$.}

\begin{Rem}[Null sets---again]
    In the following we will often make statements like:
    ``The Markov kernel $K(X_A|X_B,X_C,X_D)$ is unique up to a measurable $K(X_B|X_C,X_D)$-null set in $\Xcal_{B\cup C}$'',  (rather than in $\Xcal_{B \cup C \cup D}$). 
This means that the corresponding null set $N$ can be considered constant in $X_{D \sm (B \cup C)}$, or, more precisely, 
that $N$ is of the form:
\[ N = M \times \Xcal_{D \sm (B \cup C)}  \ins \Xcal_{B \cup C \cup D},  \]
with $M \ins \Xcal_{B\cup C}$.
\end{Rem}


\begin{Prp}[Do-calculus---existence and uniqueness]
    \label{prp:do-calculus}
    Consider an L-CBN:
    \[M= \lp G^+=\lp J,(V,U),E^+ \rp, \lp P_v ( X_v|\doit( X_{\Pa^{G^+}(v)})) \rp_{v \in V \cup U} \rp.\]
    Let $A,B,C \ins V$ and $D \ins J \cup V$ be such that $A,B,C,D$ are pairwise disjoint.\\
    \JorisDoInput{In rule 3, should we allow for $B \subseteq V \cup W \cup J$?}
    Then we have the following 3 rules relating marginal conditional to marginal interventional Markov kernels:
    \begin{enumerate}
        \item \underline{Insertion/deletion of observation}: If we have:
            \[ A \Perp^d_{G_{\doit(D)}} B \given C \cup D,   \]
            then there exists a Markov kernel:
            \[ P\lp X_A|\cancel{X_B},X_C,\doit(X_{D\cup \cancel{J}})\rp:\;  \Xcal_C \times \Xcal_{D} \dshto \Xcal_A, \]
            that is a version of:
            \[P\lp X_A|X_{B_2},X_C,\doit(X_{D\cup J})\rp,\] 
            for every subset $B_2 \ins B$ simultaneously.
            Note that this Markov kernel is only dependent on $x_C$ and $x_D$, and constant in $x_{J\setminus D}$.
        \item[] Such a Markov kernel is unique up to a measurable $P(X_C|\doit(X_{D\cup J}))$-null set in $\Xcal_{C\cup D}$. 
         \item \underline{Action/observation exchange}: If we have:
            \[ A \Perp^d_{G_{\doit(I_B,D)}} I_B \given B \cup C \cup D,   \]
            then there exists a Markov kernel:
            \[P\lp X_A|\cancel{\doit}(X_B),X_C,\doit(X_{D\cup \cancel{J}})\rp: \;\Xcal_B \times \Xcal_C \times \Xcal_{D} \dshto \Xcal_A,\]
            that is a version of:            
            \[ P\lp X_A|X_{B_1},\doit(X_{B_2}),X_C,\doit(X_{D\cup J})\rp, \]
            for every decomposition: $B=B_1\dcup B_2$, simultaneously.
        \item[] Such a Markov kernel is unique up to a measurable $P(X_B,X_C|\doit(X_{I_B},X_{D\cup J}))$-null set in $N \ins \Xcal_{B \cup C\cup D}$, i.e.\ $N$ is a $P(X_{B_1},X_C|\doit(X_{B_2},X_{D\cup J}))$-null set
            for every decomposition $B=B_1\dcup B_2$ simultaneously.
        \item \underline{Insertion/deletion of action}: If we have:
            \[ A \Perp^d_{G_{\doit(I_B,D)}} I_B \given C  \cup D,   \]
            then there exists a Markov kernel: 
            \[P\lp X_A|\cancel{\doit(X_B)},X_C,\doit(X_{D\cup \cancel{J}})\rp:\; \Xcal_C \times \Xcal_{D} \dshto \Xcal_A,\] 
            that is a version of:  
            \[ P\lp X_A|\doit(X_{B_2}),X_C,\doit(X_{D\cup J})\rp\]
            for every subset $B_2 \ins B$ simultaneously. 
            Note that this Markov kernel is only dependent on $x_C$ and $x_D$, and constant in $x_{J \sm D}$.
        \item[] Such a Markov kernel is unique up to a measurable $P\lp X_C|\doit(X_{I_B},X_{D\cup J})\rp$-null set $N \ins \Xcal_{C\cup D}$, i.e.\ $N$ is a $P(X_C|\doit(X_{B_2},X_{D\cup J}))$-null set for every subset $B_2 \ins B$ simultaneously.
    \end{enumerate}
\begin{proof}
    We make use of the global Markov property (GMP), theorem \ref{thm-gmp-mCBN}.

1.) The assumption:
            \[ A \Perp^d_{G_{\doit(D)}} B \given C \cup D,   \]
implies the conditional independence by GMP \ref{thm-gmp-mCBN}:
  \[ X_A \Indep_{P(X_V|\doit(X_{D\cup J}))} X_B \given X_C,X_D.\]
So  we get the following factorization, where we can omit the deterministic variables from $X_D$ on the left of the conditioning lines:
  \[ P\lp X_A,X_B,X_C|\doit(X_{D\cup J})\rp = Q(X_A|X_C,X_D) \otimes P\lp X_B,X_C|\doit(X_{D\cup J})\rp,\]
for some Markov kernel $Q(X_A|X_C,X_D)$. Here
$Q(X_A|X_C,X_D)$ serves as a version of the conditional Markov kernel:
  \[P\lp X_A|X_B,X_C,\doit(X_{D\cup J})\rp.\]
If we marginalize out $X_{B_1}$ for any decomposition $B=B_1 \dcup B_2$ 
in the above factorization we also get:
  \[ P\lp X_A,X_C,X_{B_2}|\doit(X_{D\cup J})\rp = Q(X_A|X_C,X_D) \otimes P\lp X_C,X_{B_2}|\doit(X_{D\cup J})\rp,\]
showing that $Q(X_A|X_C,X_D)$ is also a version of:
  \[P\lp X_A|X_{B_2},X_C,\doit(X_{D\cup J})\rp.\]
In particular, this holds for $B_2 = \emptyset$. This shows all the claimed properties for $Q(X_A|X_C,X_D)$.

Now consider another Markov kernel $K(X_A|X_C,X_D)$ and the measurable sets:
\begin{align*}
    \tilde N &:= \lC x_{C \cup D} \in \Xcal_{C \cup D} \st Q(X_A|X_C=x_C,X_D=x_D) \neq K(X_A|X_C=x_C,X_D=x_D)  \rC,\\
     N &:= \lC x_{C \cup D \cup J} \in \Xcal_{C \cup D \cup J} \st Q(X_A|X_C=x_C,X_D=x_D) \neq K(X_A|X_C=x_C,X_D=x_D)  \rC \\
       &= \tilde N \times \Xcal_{J \sm D}.
\end{align*}
If $K(X_A|X_C,X_D)$ is a version of: 
\[P(X_A|X_{B_2},X_C,\doit(X_{D\cup J})),\] 
for every subset $B_2 \ins B$ simultaneously, then this holds, in particular, for $B_2 = \emptyset$.
Since conditional Markov kernels are essentially unique, by Theorem \ref{thm-regular-conditional-Markov-kernel}, we have that $N$ is a $P(X_C|\doit(X_{D\cup J}))$-null set.
%Note that then automatically the product $\Xcal_{B_2} \times N$ is a $P(X_{B_2},X_C|\doit(X_{D\cup J}))$-null set for every $B_2 \ins B$. 
%\Joris{The last sentence proves something redundant (not part of the claim). If it is important, expand the claim accordingly. If not, then what is its purpose here?}
\\

2.) The assumption:
            \[ A \Perp^d_{G_{\doit(I_B,D)}} I_B \given B \cup C \cup D,   \]
implies the conditional independence by GMP \ref{thm-gmp-mCBN}:
  \[ X_A \Indep_{P(X_V|\doit(X_{I_B},X_{D\cup J}))} X_{I_B} \given X_B,X_C,X_D.\]
So we have the following factorization:
\begin{align}
    P\lp X_A,X_B,X_C|\doit(X_{I_B},X_{D\cup J})\rp &= Q(X_A|X_B,X_C,X_D) \otimes P\lp X_B,X_C|\doit(X_{I_B},X_{D\cup J})\rp,\label{eq:2nd-do-int-1}  
\end{align}
for some Markov kernel $Q(X_A|X_B,X_C,X_D)$, which serves as a version of the conditional Markov kernel:
  \[P\lp X_A|X_B,X_C,\doit(X_{I_B},X_{D\cup J})\rp,\]
and which is independent of $X_{I_B}$.

We fist claim that for a Markov kernel $Q(X_A|X_B,X_C,X_D)$ the equation \ref{eq:2nd-do-int-1} is equivalent to the system of equations \ref{eq:2nd-do-int-2} indexed by subsets $B_2 \ins B$ and with $B_1 := B \sm B_2$:
\begin{align}
    P\lp X_A,X_{B_1},X_C|\doit(X_{B_2},X_{D\cup J})\rp &= Q(X_A|X_B,X_C,X_D) \otimes P\lp X_{B_1},X_C|\doit(X_{B_2},X_{D\cup J})\rp. \label{eq:2nd-do-int-2} 
\end{align}
Indeed, we can look at the different input values for
$X_{I_{B}} = (X_{I_{B_1}},X_{I_{B_2}})$ in equation \ref{eq:2nd-do-int-1}.
For $B_1$ we put: $X_{I_{B_1}}= \star = (\star)_{v \in B_1}$ and for $B_2$ we take values: 
$X_{I_{B_2}}=x_{B_2} \in \Xcal_{B_2}$.
This implies: 
\begin{align*} 
  & P\lp X_A,X_B,X_C|\doit(X_{{B_2}}=x_{B_2},X_{D\cup J})\rp \\
  &= P\lp X_A,X_B,X_C|\doit(X_{I_{B_1}}=\star,X_{I_{B_2}}=x_{B_2},X_{D\cup J})\rp \\
  &= Q(X_A|X_B,X_C,X_D) \otimes P\lp X_B,X_C|\doit(X_{I_{B_1}}=\star,X_{I_{B_2}}=x_{B_2},X_{D\cup J})\rp,\\
  &= Q(X_A|X_B,X_C,X_D) \otimes P\lp X_B,X_C|\doit(X_{{B_2}}=x_{B_2},X_{D\cup J})\rp.
\end{align*}
So we get the equations:
\begin{align*}
    P\lp X_A,X_B,X_C|\doit(X_{B_2},X_{D\cup J})\rp 
  &= Q(X_A|X_B,X_C,X_D) \otimes P\lp X_B,X_C|\doit(X_{B_2},X_{D\cup J})\rp, 
\end{align*}
  where we can further marginalize out the deterministic $X_{B_2}$:
  \begin{align*} P\lp X_A,X_{B_1},X_C|\doit(X_{B_2},X_{D\cup J})\rp 
  &= Q(X_A|X_B,X_C,X_D) \otimes P\lp X_{B_1},X_C|\doit(X_{B_2},X_{D\cup J})\rp. %\label{eq:2nd-do-int-3} 
  \end{align*}
  Note  that we can also go back by multiplying with $\delta(X_{B_2}|X_{B_2})$. This shows the intermediate claim.

The equation \ref{eq:2nd-do-int-2} already implies that $Q(X_A|X_B,X_C,X_D)$ is a version of the conditional Markov kernel:
\begin{align} P\lp X_A|X_{B_1},X_C,\doit(X_{B_2},X_{D\cup J})\rp, \label{eq:2nd-do-int-4} 
\end{align}
for every decomposition: $B=B_1 \dcup B_2$ simultaneously.

Now consider another Markov kernel $K(X_A|X_B,X_C,X_D)$ and the measurable sets:
\begin{align*}
    \tilde N &:= \bigl\{ x_{B \cup C \cup D} \in \Xcal_{B \cup C \cup D} \,\big|\,  Q(X_A|X_B=x_B,X_C=x_C,X_D=x_D) \\ 
    & \qquad \qquad  \neq K(X_A|X_B=x_B,X_C=x_C,X_D=x_D)  \bigr\},\\
  N &:= \tilde N \times \Xcal_{J \sm D}.
\end{align*}
If $K(X_A|X_B,X_C,X_D)$ now is also a version of the conditional Markov kernel \ref{eq:2nd-do-int-4} 
for every decomposition $B=B_1 \dcup B_2$ simultaneously, then $N$ is a 
$P(X_{B_1},X_C|\doit(X_{B_2},X_{D\cup J}))$-null set for every decomposition $B=B_1 \dcup B_2$, 
because conditional Markov kernels are essentially unique, see Theorem \ref{thm-regular-conditional-Markov-kernel}. 
By Lemma \ref{lem:int-null-set} this statement is equivalent for $N$ to be a 
$P(X_B,X_C|\doit(X_{I_B},X_{D\cup J}))$-null set.\\


3.) The assumption:
            \[ A \Perp^d_{G_{\doit(I_B,D)}} I_B \given C  \cup D,   \]
implies the conditional independence by GMP \ref{thm-gmp-mCBN}:
\[ X_A \Indep_{P(X_V|\doit(X_{I_B},X_{D\cup J}))} X_{I_B} \given X_C,X_D.\]
So we have the following factorization:
\begin{align} 
    P\lp X_A,X_C|\doit(X_{I_B},X_{D\cup J})\rp &= Q(X_A|X_C,X_D) \otimes P\lp X_C|\doit(X_{I_B},X_{D\cup J})\rp,
\label{eq:3nd-do-int-1} 
\end{align}
for some Markov kernel $Q(X_A|X_C,X_D)$, which serves as a version of the conditional Markov kernel:
\[P\lp X_A|X_C,\doit(X_{I_B},X_{D\cup J})\rp,\]
and which is independent of $X_{I_B}$.\\ 
We can now look at the different input values for any decomposition: $B=B_1 \dcup B_2$.
For this we put: $X_{I_{B_1}}= \star = (\star)_{v \in B_1}$ and  $ X_{I_{B_2}}=x_{B_2} \in \Xcal_{B_2}$.
This implies:
\begin{align*} 
  & P\lp X_A,X_C|\doit(X_{{B_2}}=x_{B_2},X_{D\cup J})\rp \\
  &= P\lp X_A,X_C|\doit(X_{I_{B_1}}=\star,X_{I_{B_2}}=x_{B_2},X_{D\cup J})\rp \\
  &= Q(X_A|X_C,X_D) \otimes P\lp X_C|\doit(X_{I_{B_1}}=\star,X_{I_{B_2}}=x_{B_2},X_{D\cup J})\rp,\\
  &= Q(X_A|X_C,X_D) \otimes P\lp X_C|\doit(X_{{B_2}}=x_{B_2},X_{D\cup J})\rp.
\end{align*}
So we get for every subset $B_2 \ins B$:
\begin{align} 
    P\lp X_A,X_C|\doit(X_{B_2},X_{D\cup J})\rp &= Q(X_A|X_C,X_D) \otimes P\lp X_C|\doit(X_{B_2},X_{D\cup J})\rp, \label{eq:3nd-do-int-2} 
\end{align}
which shows that $Q(X_A|X_C,X_D)$ is a version of the conditional Markov kernel:
\begin{align} P\lp X_A|X_C,\doit(X_{B_2},X_{D\cup J})\rp, \label{eq:3nd-do-int-3}  \end{align}
for every subset $B_2 \ins B$ simultaneously.

Now consider another Markov kernel $K(X_A|X_C,X_D)$ and the measurable sets:
\begin{align*}
    \tilde N &:= \lC x_{C \cup D} \in \Xcal_{C \cup D} \st Q(X_A|X_C=x_C,X_D=x_D) \neq K(X_A|X_C=x_C,X_D=x_D)  \rC,\\
    N_{B_2} &:= \Xcal_{B_2} \times \tilde N \times \Xcal_{J \sm D}.
\end{align*}
Now assume that $K(X_A|X_C,X_D)$ is a version of the conditional Markov kernel in (\ref{eq:3nd-do-int-3})
for every subset $B_2 \ins B$ simultaneously. 
Then for every subset $B_2 \ins B$ set $N_{B_2}$ is a
$P(X_C|\doit(X_{B_2},X_{D\cup J}))$-null set, because of the essential uniqueness of conditional Markov kernels, see Theorem \ref{thm-regular-conditional-Markov-kernel}. More concretely, for $x_{B_2} \in \Xcal_{B_2}$ 
and $x_{D \cup J} \in \Xcal_{D \cup J}$ we get the equations:
\begin{align*}
  0 &=P(X_C \in (N_{B_2})_{(x_{B_2},x_{D\cup J})} |\doit(X_{B_2}=x_{B_2},X_{D\cup J}=x_{D \cup J})) \\
      &=P(X_C \in \tilde N_{x_D} |\doit(X_{B_2}=x_{B_2},X_{D\cup J}=x_{D \cup J}))\\
      &=P(X_C \in \tilde N_{x_D} |\doit(X_{I_{B_1}}=\star, X_{I_{B_2}}=x_{B_2},X_{D\cup J}=x_{D \cup J}))\\
      &=P(X_C \in \tilde N_{x_D} |\doit(X_{I_B}=(\star,x_{B_2}),X_{D\cup J}=x_{D \cup J})).
\end{align*}
Since we have this for all decompositions $B=B_1 \dcup B_2$ and all values $x_{B_2} \in \Xcal_{B_2}$ 
we are running through all values $x_{I_B} \in \Xcal_{I_B}$. 
This shows that $\tilde N$ is a $P(X_C|\doit(X_{I_B},X_{D\cup J}))$-null set in $\Xcal_{C \cup D}$.
\end{proof}
\end{Prp}

\restatethmasdocalculusdetail*
\begin{proof}
 W.l.o.g.\ we can assume all $\mu_v$ to be probability measures.

 1.) Let $Q(X_A|X_C,X_D)$ be the Markov kernel from Proposition \ref{prp:do-calculus} point 1.
Recall that conditional Markov kernels are essentially unique by Theorem \ref{thm-regular-conditional-Markov-kernel}. 
    This shows that the set:
\begin{align*}
    \tilde N^{(i)} := \bigl\{ x_{B^{(i)} \cup C \cup D \cup J} \in \Xcal_{B^{(i)} \cup C \cup D \cup J} \,\big|\,
       & Q(X_A|X_C=x_C,X_D=x_D)\\ 
       &\neq P(X_A|X_{B^{(i)}}=x_{B^{(i)}},X_C=x_C,\doit(X_{D\cup J}=x_{D \cup J})) \bigr\},
\end{align*}
is a (measurable) $P(X_{B^{(i)}},X_C|\doit(X_{D\cup J}))$-null set. So the lifted set:
\begin{align*}
    N^{(i)} := \bigl\{ x_{B \cup C \cup D \cup J} \in \Xcal_{B \cup C \cup D \cup J} \,\big|\,
       & Q(X_A|X_C=x_C,X_D=x_D)\\ 
       &\neq P(X_A|X_{B^{(i)}}=x_{B^{(i)}},X_C=x_C,\doit(X_{D\cup J}=x_{D \cup J})) \bigr\},
\end{align*}
is then a (measurable) $P(X_B,X_C|\doit(X_{D\cup J}))$-null set. Then also the finite union:
 \[ N := \bigcup_{i \in I} N^{(i)}  \ins \Xcal_{B \cup C \cup D \cup J},  \]
 is a (measurable) $P(X_B,X_C|\doit(X_{D\cup J}))$-null set as well. 
 Note that on the complement $N^\cmpl$ all Markov kernels agree with $Q(X_A|X_C,X_D)$ and are thus all equal on $N^\cmpl$.\\

2.) Consider the Markov kernel $Q(X_A|X_B,X_C,X_D)$ from Proposition \ref{prp:do-calculus} point 2
    and for $i \in I$ the measurable set:
\begin{align*}
     N^{(i)} :=  \bigl\{ & x_{B \cup C \cup D \cup J} \in \Xcal_{B \cup C \cup D \cup J} \,\big|\,
        Q(X_A|X_B=x_B,X_C=x_C,X_D=x_D)\\ 
       &\neq P(X_A|X_{B_1^{(i)}}=x_{B_1^{(i)}},\doit(X_{B_2^{(i)}}=x_{B_2^{(i)}}),X_C=x_C,\doit(X_{D\cup J}=x_{D\cup J}))  \bigr\}.
\end{align*}
Again, by the essential uniqueness of conditional Markov kernels, Theorem \ref{thm-regular-conditional-Markov-kernel}, the set $N^{(i)}$ is 
%a $P(X_{B_1^{(i)}},X_C|\doit(X_{B_2^{(i)}},X_{D\cup J}))$-null set and thus
a $P(X_{B_1^{(i)}},X_C|\doit(X_{B_2^{(i)}},X_{D\cup J})) \otimes \mu_{B_2^{(i)}}$-null set.
The absolute continuity:
 \[ \mu_{B \cup C} \ll  P(X_{B_1^{(i)}},X_C|\doit(X_{B_2^{(i)}},X_{D\cup J})) \otimes \mu_{B_2^{(i)}}, \]
 then renders $N^{(i)}$ a $\mu_{B \cup C}$-null set. 
 This shows that the finite union:
 \[ N := \bigcup_{i \in I} N^{(i)}, \]
 is a $\mu_{B \cup C}$-null set as well.  Again, note that on the complement $N^\cmpl$ all Markov kernels agree with $Q(X_A|X_B,X_C,X_D)$ and are thus all equal on $N^\cmpl$. \\

3.) Consider the Markov kernel $Q(X_A|X_C,X_D)$ from Proposition \ref{prp:do-calculus} point 3 and
for $i \in I$ the measurable set:
\begin{align*}
    \tilde N^{(i)} := \bigl\{ x_{B^{(i)} \cup C \cup D \cup J} &\in \Xcal_{B^{(i)} \cup C \cup D \cup J} \,\big|\,
        Q(X_A|X_C=x_C,X_D=x_D)\\ 
       &\neq P(X_A|\doit(X_{B^{(i)}}=x_{B^{(i)}}),X_C=x_C,\doit(X_{D\cup J}=x_{D \cup J})) \bigr\}.
\end{align*}
Again, by the essential uniqueness of conditional Markov kernels, Theorem \ref{thm-regular-conditional-Markov-kernel}, the set $\tilde N^{(i)}$ is a $P(X_C|\doit(X_{B^{(i)}},X_{D\cup J}))$-null set. By the absolute continuity $ \mu_C \ll P(X_C|\doit(X_{B^{(i)}},X_{D\cup J}))$ we get that $\tilde N^{(i)}$ is a $\mu_C$-null set. This shows that the measurable set:
\begin{align*}
    N^{(i)} := \bigl\{ x_{B \cup C \cup D \cup J} &\in \Xcal_{B \cup C \cup D \cup J} \,\big|\,
        Q(X_A|X_C=x_C,X_D=x_D)\\ 
       &\neq P(X_A|\doit(X_{B^{(i)}}=x_{B^{(i)}}),X_C=x_C,\doit(X_{D\cup J}=x_{D \cup J})) \bigr\}.
\end{align*}
is a $\mu_C$-null set as well. Then the finite union:
 \[ N := \bigcup_{i \in I} N^{(i)}, \]
 is also a $\mu_C$-null set. 
Again, note that on the complement $N^\cmpl$ all Markov kernels agree with $Q(X_A|X_C,X_D)$ and are thus all equal on $N^\cmpl$.
\end{proof}
